\section{Síntesis y ampliación}

\subsection{Puntos clave}

Esta clase parte de VAE y deriva progresivamente el marco básico de los modelos de difusión. Los puntos clave incluyen:

\begin{enumerate}
    \item \textbf{Limitaciones de VAE}: modelar la verosimilitud y la posterior como gaussianas es una simplificación excesiva, lo que limita la calidad generada.
    \item \textbf{Enfoque jerárquico}: introduce variables latentes multinivel; cada nivel sigue siendo una gaussiana simple, pero la distribución global es más compleja.
    \item \textbf{Suposición de Markov}: simplifica el cálculo del VAE jerárquico, permitiendo una descomposición eficiente del ELBO.
    \item \textbf{Diseño del proceso hacia adelante}: agrega ruido gradualmente mediante un esquema de ruido predefinido $\{\beta_t\}$, lo cual tiene solución en forma cerrada.
    \item \textbf{Aprendizaje del proceso inverso}: solo es necesario aprender el proceso de eliminación de ruido $p_\theta(x_{t-1}|x_t)$, evitando la dificultad del entrenamiento conjunto.
    \item \textbf{Descomposición del ELBO}: el término de ajuste al prior se satisface automáticamente; el entrenamiento se centra en el término de consistencia.
\end{enumerate}

\subsection{Lecturas adicionales}

\begin{itemize}
    \item Ho, J., Jain, A., \& Abbeel, P. (2020). \textit{Denoising Diffusion Probabilistic Models}. NeurIPS 2020. --- Artículo original de DDPM.
    \item Sohl-Dickstein, J., et al. (2015). \textit{Deep Unsupervised Learning using Nonequilibrium Thermodynamics}. ICML 2015. --- Artículo teórico que fundamenta los modelos de difusión.
    \item Kingma, D. P. \& Welling, M. (2014). \textit{Auto-Encoding Variational Bayes}. ICLR 2014. --- Artículo original de VAE.
    \item Luo, C. (2022). \textit{Understanding Diffusion Models: A Unified Perspective}. arXiv:2208.11970. --- Revisión desde una perspectiva unificada.
    \item Nichol, A. Q. \& Dhariwal, P. (2021). \textit{Improved Denoising Diffusion Probabilistic Models}. ICML 2021. --- Mejoras sobre esquemas de ruido coseno y otros aspectos.
\end{itemize}

%% --- Fin del contenido --- %%

\end{document}
